% Options for packages loaded elsewhere
\PassOptionsToPackage{unicode}{hyperref}
\PassOptionsToPackage{hyphens}{url}
%
\documentclass[
]{article}
\usepackage{amsmath,amssymb}
\usepackage{lmodern}
\usepackage{ifxetex,ifluatex}
\ifnum 0\ifxetex 1\fi\ifluatex 1\fi=0 % if pdftex
  \usepackage[T1]{fontenc}
  \usepackage[utf8]{inputenc}
  \usepackage{textcomp} % provide euro and other symbols
\else % if luatex or xetex
  \usepackage{unicode-math}
  \defaultfontfeatures{Scale=MatchLowercase}
  \defaultfontfeatures[\rmfamily]{Ligatures=TeX,Scale=1}
\fi
% Use upquote if available, for straight quotes in verbatim environments
\IfFileExists{upquote.sty}{\usepackage{upquote}}{}
\IfFileExists{microtype.sty}{% use microtype if available
  \usepackage[]{microtype}
  \UseMicrotypeSet[protrusion]{basicmath} % disable protrusion for tt fonts
}{}
\makeatletter
\@ifundefined{KOMAClassName}{% if non-KOMA class
  \IfFileExists{parskip.sty}{%
    \usepackage{parskip}
  }{% else
    \setlength{\parindent}{0pt}
    \setlength{\parskip}{6pt plus 2pt minus 1pt}}
}{% if KOMA class
  \KOMAoptions{parskip=half}}
\makeatother
\usepackage{xcolor}
\IfFileExists{xurl.sty}{\usepackage{xurl}}{} % add URL line breaks if available
\IfFileExists{bookmark.sty}{\usepackage{bookmark}}{\usepackage{hyperref}}
\hypersetup{
  pdftitle={Educational Polarization in the CES},
  hidelinks,
  pdfcreator={LaTeX via pandoc}}
\urlstyle{same} % disable monospaced font for URLs
\usepackage[margin=1in]{geometry}
\usepackage{graphicx}
\makeatletter
\def\maxwidth{\ifdim\Gin@nat@width>\linewidth\linewidth\else\Gin@nat@width\fi}
\def\maxheight{\ifdim\Gin@nat@height>\textheight\textheight\else\Gin@nat@height\fi}
\makeatother
% Scale images if necessary, so that they will not overflow the page
% margins by default, and it is still possible to overwrite the defaults
% using explicit options in \includegraphics[width, height, ...]{}
\setkeys{Gin}{width=\maxwidth,height=\maxheight,keepaspectratio}
% Set default figure placement to htbp
\makeatletter
\def\fps@figure{htbp}
\makeatother
\setlength{\emergencystretch}{3em} % prevent overfull lines
\providecommand{\tightlist}{%
  \setlength{\itemsep}{0pt}\setlength{\parskip}{0pt}}
\setcounter{secnumdepth}{-\maxdimen} % remove section numbering
\ifluatex
  \usepackage{selnolig}  % disable illegal ligatures
\fi

\title{Educational Polarization in the CES}
\author{}
\date{\vspace{-2.5em}2021-11-23}

\begin{document}
\maketitle

Recently, CES Researcher \href{https://twitter.com/DeshpandePia}{Pia
Deshpande} wrote an excellent tutorial detailing how to
\href{https://pdeshlab.github.io/posts/CES-plots-over-time}{plot trends
over time} using CES data --- I highly encourage anyone who hasn't yet
read the piece to promptly do so! This style of plot pairs particularly
well with the
\href{https://dataverse.harvard.edu/dataset.xhtml?persistentId=doi:10.7910/DVN/II2DB6}{cumulative
CES Common Content} dataset created by
\href{https://www.shirokuriwaki.com/}{Shiro Kuriwaki}. Together, these
resources can help us to understand an issue that has received
\href{https://www.nytimes.com/2021/09/08/us/politics/how-college-graduates-vote.html}{increased
attention} in recent months: educational polarization.

Research has noted
\href{https://www.pewresearch.org/politics/2018/03/20/1-trends-in-party-affiliation-among-demographic-groups/}{an
increasing divide} in party identification along educational lines. But
less attention has been devoted to how these changes differ by race.
Using CES data, I note that across most racial groups, Americans with a
college education less strongly identify as Republicans while Americans
without a college degree less strongly identify as Democrats. These
trends in partisan identification are complex, however, and differ in
strength and effect by race. When it comes to vote choice, however,
educational polarization is mostly limited to white Americans.

\hypertarget{americans-without-a-college-education-of-all-races-are-increasingly-unlikely-to-identify-as-democrats}{%
\section{Americans without a college education of all races are
increasingly unlikely to identify as
Democrats}\label{americans-without-a-college-education-of-all-races-are-increasingly-unlikely-to-identify-as-democrats}}

\includegraphics{/images/posts/educational-polarization/pid-dem-race.png}

The percent of those without a college degree identifying as Democrats
sharply declined between 2012 and 2020, even as the share of Americans
with a college degree identifying as Democrats remained somewhat
constant between 2008 and 2020. While educational polarization in vote
choice was mostly limited to nonhispanic white Americans, these trends
in party identification played out in remarkably similar ways among
Americans of all races.

Among both Black Americans\footnote{It is worth noting that Black
  Americans with a college degree were the only college-educated racial
  group to identify as Democrats at a statistically significant lower
  rate in 2020 than in 2016 and 2012; this decline, however, was less
  pronounced than the attrition among those without a college degree so
  still led to a statistically significant educational gap.} and
Hispanic Americans, there seemed to be no educational differences in the
share of self-identified Democrats in 2008 and 2012; by 2020, however,
as Democratic self-identification declined among those without a college
education, relatively narrow --- yet statistically significant --- gaps
of about eight percentage points emerged.

In 2008 and 2012, there were more clear differences along educational
lines in the share of Asian and white Americans identifying as
Democrats.\footnote{Like before, the relatively smaller sample size of
  Asian Americans in the CES leads to large errors for the point
  estimates, so it is an important caveat to note that the educational
  differences in 2008 and 2012 are not statistically significant. Still,
  following the overall trends these differences grow and by 2016 and
  2020 the difference in the point estimates are statistically
  significant.} These divides only widened in 2016 and by 2020,
college-educated white and Asian Americans identified as Democrats at a
rate almost 20 percentage points higher than their non-college educated
peers.

\hypertarget{only-white-americans-without-a-college-degree-were-sharply-more-likely-to-identify-as-republicans}{%
\section{Only white Americans without a college degree were sharply more
likely to identify as
Republicans}\label{only-white-americans-without-a-college-degree-were-sharply-more-likely-to-identify-as-republicans}}

There is strong evidence that Americans of all races without a college
degree are increasingly less likely to self-identify as Democrats. How
has education level impacted partisan identification specifically, then,
for Americans of different races? Is the Democratic attrition among
Americans without a college education being driven by switchers to the
Republican party, or just an increase in independents?

\includegraphics{/images/posts/educational-polarization/pid-rep-race.png}

The share of self-identified Republicans with a college degree declined
modestly between 2008 and 2020 among white, Hispanic and Asian
Americans. In contrast, though just 10 percent did, Black Americans with
a college degree were actually slightly more likely to identify as
Republican than in previous years.

The share of Hispanic and Asian Americans without a college education
who self-identify as Republicans appears to have somewhat increased,
though these differences are not statistically significant. There was a
statistically significant increase among Black Americans without a
college degree, though just 9 percent self-identified as Republicans.

Among white Americans without a college degree, however, the differences
were stark. The share of white Americans without a college degree who
identify as Republicans has grown from 40 percent in 2008 to 53 percent
in 2020.

\hypertarget{black-americans-without-a-college-degree-increasingly-identify-as-independents}{%
\section{Black Americans without a college degree increasingly identify
as
Independents}\label{black-americans-without-a-college-degree-increasingly-identify-as-independents}}

The attrition in the share of self-identified Republicans among white,
Hispanic and Asian Americans with a college degree seems to be driven
somewhat by a switch to identfying as Independents; these groups all
increasingly identified as Independents between 2008 and 2020.

Among those without a college education, the trends were less uniform.
The rate of self-identifying as an Independent appears generally stable
among white and Hispanic Americans. The estimate for the share of Asian
Americans without a college degree who identified as Independents seems
to have increased in 2020 from 2008, but due to large standard errors
these differences should not be seen as statistically significant.

\includegraphics{/images/posts/educational-polarization/pid-ind-race.png}

Black Americans without a college education saw the most notable
increase in the rate at which they identified as Independents. The share
of self-identified Independent Black Americans without a college degree
in 2020 was double the 2008 level, though at just 16 percent this still
marked a small share overall.

\hypertarget{educational-polarization-in-vote-choice-is-largely-driven-by-nonhispanic-whites}{%
\section{Educational polarization in vote choice is largely driven by
nonhispanic
whites}\label{educational-polarization-in-vote-choice-is-largely-driven-by-nonhispanic-whites}}

But have these shifts in party identification led to educational
polarization in vote choice? Examining support for the Democratic
presidential candidate by race and education between 2008 and 2020
reveals a few notable trends. There appears to have been little
variation in Hispanic support for the Democratic presidential candidate
by education. While these levels appear to have declined slightly from
their 2012 levels, we cannot conclude that the difference is
statistically significant.

\includegraphics{/images/posts/educational-polarization/voted-race.png}

We see a slight but statistically significant drop in Democratic support
among Black Americans between 2012 and 2020. For Black Americans there
appears to be no discernible difference, however, along educational
lines, and Democratic support in 2020 appears to have remained at a
comparable level to 2016.

Estimates for Democratic presidential support among Asian Americans
appear noisy, in part due to smaller sample sizes. This makes it
difficult to draw strong conclusions about trends between groups and
election years. In 2020, however, CES data noted that Asian Americans
with a college education voted for Joe Biden at a statistically
significant higher rate than Asian Americans without a college
education.

It is hardly
\href{https://www.washingtonpost.com/news/monkey-cage/wp/2014/01/23/how-better-educated-whites-are-driving-political-polarization/}{a
novel observation}, but nonhispanic white Americans appear to most
strongly drive educational polarization. In 2008 and 2012 white
Americans with a college education voted for the Democratic candidate at
a rate just nine percentage points higher than white Americans without a
college education; this gap more than doubled by 2020, with more than
half of white Americans with a college education voting for Joe Biden,
compared to just 30 percent of those without a college degree.

\hypertarget{conclusion}{%
\section{Conclusion}\label{conclusion}}

There appear to be growing educational divides in party identification
among Americans of all races. Hispanic, white, and Asian Americans with
a college degree have decreasingly identified as Republicans and
increasingly identified as Independents.

Americans without a college degree of all races identify less as
Democrats but have sorted in different ways. Black Americans without a
college degree identify slightly more as Republicans but more strongly
as Independents than in previous years, while white Americans without a
college degree are sharply more likely to identify as Republicans.

Despite these trends, educational polarization in vote choice is mostly
limited to white Americans but may be emerging among Asian Americans.
Why has vote choice remained relatively stable in the face of shifting
partisan identification?

The absolute magnitude of these trends are somewhat small, and
asymmetric: even as the share of Black Americans without a college
degree who identify as Democrats dropped by 15 percentage points, the
share who identified as Republicans rose by just 3 percentage points.
Partisans are not defecting to the other party, but become increasingly
likely to identify as Independents or respond ``Not Sure''. And while
education seems to moderate partisan realignment, only white Americans
differ in vote choice along educational lines.

An excellent
\href{https://www.dropbox.com/s/3nsueuaoquxzki1/realignments_draft.pdf?dl=0}{working
paper} by \href{http://michaeljaybarber.com/}{Michael Barber} and
\href{http://jeremypope.byu.edu/}{Jeremy C. Pope} --- published while I
was drafting this blog post --- presents a more thorough discussion of
the dynamics and implications of these political shifts. As they note,
further research on realignments must account for the potential
interaction and confounding effects of race.

\end{document}
